\section{Properties observations and assurance boundaries}\label{sec:C25:properties}
Start with a property $P(x,e)$ of artifact or action $x$ in environment $e$. A verifier $V$ observes a projection $o(x,e)$ and returns accept, reject, or unresolved. The projection can omit decisive state. For example, a screenshot may not show a file's location, and a test log may not identify the artifact revision. No verifier can reliably recover distinctions absent from its observation without additional assumptions or queries.

This leads to three separate questions. Is the property well defined? Does the observation contain evidence sufficient to judge it? Does the verifier use that evidence correctly? A weak result at the third stage can be improved by a stronger judge. A failure at the second stage requires a better observation interface. A failure at the first stage requires specification or adjudication, not merely more inference. Collapsing these cases into one accuracy score obscures the intervention needed.

Formal verification supplies strong conclusions when the model, specification, and proof obligations are appropriate. Its guarantee remains conditional on those objects representing the deployed system. A proof about a pure function does not establish the behavior of a surrounding agent that calls another version, supplies invalid inputs, or operates with changed permissions. The deployment boundary must bind the verified artifact and assumptions to actual execution.

Executable tests provide finite observations. They can falsify a candidate under the tested conditions and can estimate behavior over a sampled distribution. Passing them does not prove a universal property. Nonetheless, a reproducible failing test is stronger evidence than an unsupported critique because another investigator can replay the witness. The key is to preserve the relationship between the witness and the intended requirement rather than accept every disagreement with a reference implementation as a defect.

An assurance case organizes these distinctions. Its claims identify required properties; evidence records observations and proofs; arguments explain why the evidence supports the claims; defeaters identify conditions that invalidate the argument. This structure is especially valuable for agents because the system changes across sessions, memory updates, model releases, and tool schemas. Assurance must track those dependencies instead of treating one successful evaluation as permanent certification.

For a binary property $Y$ and a fixed observation $O$, let $\eta(o)=\Pr(Y=1\mid O=o)$. Under zero--one loss, the minimum error over all measurable classifiers using only $O$ is
\begin{equation}
 R_O^*=\mathbb E\big[\min\{\eta(O),1-\eta(O)\}\big].
\end{equation}
At each observation, choosing the more probable class minimizes conditional error; averaging gives the expression. It is a population identity, not an estimate available without labeled data. Two equally probable states with identical observations and opposite property values contribute an irreducible error of one half on that ambiguous subset. More inference over the same evidence cannot remove the ambiguity. Additional authoritative queries can change $O$ and therefore change the bound; a richer explanation generated from the original $O$ cannot add independent evidence. This supplies a formal reason to diagnose observation loss separately from judge weakness.
